% Original timeline; winding path inspired by user reference, with no decorative icons.
\begin{tikzpicture}[x=1mm,y=1mm,font=\figurefont\footnotesize,
 milestone/.style={circle,minimum size=7mm,inner sep=0pt,draw=BookTeal,fill=BookTeal!20,line width=.8pt},
 title/.style={font=\figurefont\small,text=BookInk},
 detail/.style={font=\figurefont\footnotesize,text=BookMuted}]
\path[use as bounding box] (0,-8) rectangle (155,182);
\draw[BookMuted!60,line width=3.5mm,line cap=round]
 (78,178) .. controls (48,174) and (48,164) .. (78,158)
 .. controls (108,152) and (108,142) .. (78,136)
 .. controls (48,130) and (48,120) .. (78,114)
 .. controls (108,108) and (108,98) .. (78,92)
 .. controls (48,86) and (48,76) .. (78,70)
 .. controls (108,64) and (108,54) .. (78,48)
 .. controls (48,42) and (48,32) .. (78,26)
 .. controls (108,20) and (108,10) .. (78,4);
\draw[white,dashed,line width=.5mm,dash pattern=on 2mm off 2mm]
 (78,178) .. controls (48,174) and (48,164) .. (78,158)
 .. controls (108,152) and (108,142) .. (78,136)
 .. controls (48,130) and (48,120) .. (78,114)
 .. controls (108,108) and (108,98) .. (78,92)
 .. controls (48,86) and (48,76) .. (78,70)
 .. controls (108,64) and (108,54) .. (78,48)
 .. controls (48,42) and (48,32) .. (78,26)
 .. controls (108,20) and (108,10) .. (78,4);
\node[milestone] at (56,170) {};
\node[anchor=west,text=BookTeal,font=\figurefont\small] at (2,175) {1925};
\node[title,anchor=west,text width=49mm,align=left] at (2,169) {Ising模型};
\node[detail,anchor=north west,text width=49mm,align=left] at (2,165) {以局部相互作用描述整体配置};
\node[milestone] at (100,148) {};
\node[anchor=west,text=BookTeal,font=\figurefont\small] at (109,153) {1977};
\node[title,anchor=west,text width=46mm,align=left] at (109,147) {EM方法};
\node[detail,anchor=north west,text width=46mm,align=left] at (109,143) {推断潜变量后更新模型参数};
\node[milestone] at (56,126) {};
\node[anchor=west,text=BookTeal,font=\figurefont\small] at (2,131) {1988};
\node[title,anchor=west,text width=49mm,align=left] at (2,125) {概率网络与团树};
\node[detail,anchor=north west,text width=49mm,align=left] at (2,121) {连接独立性语义与局部概率计算};
\node[milestone] at (100,104) {};
\node[anchor=west,text=BookTeal,font=\figurefont\small] at (109,109) {1996};
\node[title,anchor=west,text width=46mm,align=left] at (109,103) {图模型理论};
\node[detail,anchor=north west,text width=46mm,align=left] at (109,99) {系统整理图、分离与分解关系};
\node[milestone] at (56,82) {};
\node[anchor=west,text=BookTeal,font=\figurefont\small] at (2,87) {2001};
\node[title,anchor=west,text width=49mm,align=left] at (2,81) {因子图};
\node[detail,anchor=north west,text width=49mm,align=left] at (2,77) {显式组织变量与局部函数};
\node[milestone] at (100,60) {};
\node[anchor=west,text=BookTeal,font=\figurefont\small] at (109,65) {2008};
\node[title,anchor=west,text width=46mm,align=left] at (109,59) {变分推断框架};
\node[detail,anchor=north west,text width=46mm,align=left] at (109,55) {以优化近似难以计算的后验};
\node[milestone] at (56,38) {};
\node[anchor=west,text=BookTeal,font=\figurefont\small] at (2,43) {2013—2015};
\node[title,anchor=west,text width=49mm,align=left] at (2,37) {神经生成模型};
\node[detail,anchor=north west,text width=49mm,align=left] at (2,33) {发展VAE、可逆流与扩散路径};
\node[milestone] at (100,16) {};
\node[anchor=west,text=BookTeal,font=\figurefont\small] at (109,21) {2020年代};
\node[title,anchor=west,text width=46mm,align=left] at (109,15) {条件生成与可控推断};
\node[detail,anchor=north west,text width=46mm,align=left] at (109,11) {兼顾分布表达、学习与采样成本};
\end{tikzpicture}
